\documentclass[11pt]{article}
\usepackage[margin=25mm]{geometry}
\usepackage{fontspec}
\setmainfont{Latin Modern Roman}
\usepackage{amsmath,amssymb,mathtools}
\usepackage{unicode-math}
\setmathfont{Latin Modern Math}
\usepackage{xurl,hyperref}
\hypersetup{colorlinks=true,urlcolor=blue}
\setcounter{secnumdepth}{0}
\setlength{\parindent}{0pt}
\setlength{\parskip}{5pt}
\setlength{\emergencystretch}{3em}
\newcommand{\sref}[2]{\hyperlink{source:#1:#2}{[S#2]}}
\newcommand{\cataloguescope}[1]{\paragraph{Recorded target.}#1}
\newcommand{\R}{\mathbb R}
\newcommand{\C}{\mathbb C}
\newcommand{\N}{\mathbb N}
\newcommand{\Z}{\mathbb Z}
\newcommand{\Q}{\mathbb Q}
\newcommand{\F}{\mathbb F}
\newcommand{\eps}{\varepsilon}
\begin{document}
% UnsolvedMath ID 30002839; source catalogue code OWR-13673-008
\setcounter{section}{30002838}
\section{Local Conformal Field Theories from Non-Semisimple Tensor Categories}\label{30002839}
{\small UnsolvedMath ID 30002839 · Mathematical Physics}
\subsection{Definitions and mathematical statement}

\paragraph{Shared notation.}$\mathbb N=\{1,2,\ldots\}$; the standard number fields have their usual meanings. All additional conventions are stated locally.


\paragraph{Definitions and conventions.}
A finite tensor category is a rigid abelian tensor category with finite-length and finite-dimensionality properties; nonsemisimple means objects need not split into simple summands. Chiral data describe one holomorphic sector of a two-dimensional conformal theory. A bulk object combines the two chiral sectors. Sewing compatibility requires correlators to agree under gluing surfaces; mapping-class-group invariance requires compatibility with the topological symmetries of those surfaces.

\paragraph{Target retained from the canonical source catalogue.}
\paragraph{Catalogue wording retained for traceability.}
For a nonsemisimple finite tensor category $\mathcal C$, do full local conformal field theories with chiral data $\mathcal C$ exist, and how can their bulk objects and sewing-compatible, mapping-class-group-invariant correlators be constructed?

\paragraph{Source scope and qualification.}
Braiding, duality and any nondegeneracy hypotheses needed by a construction must be read from the report. The extraction\textquotesingle s bare finite-tensor-category hypothesis is not silently supplemented to claim a theorem for a narrower class.
\subsection{Short English statement}
Construct full local conformal theories from nonsemisimple chiral tensor categories, including their bulk objects and consistent correlators.
\subsection{Research attempt}
\paragraph{Model, date and investigated scope.}
OpenAI GPT-6 Astra (configured model \texttt{gpt-6-astra}), ultra reasoning; generated 11 October 2026. Computes a nonsemisimple Frobenius pairing and elementary sewing contractions, exposing the separation between algebraic consistency and Hilbert positivity.

\paragraph{Sources and literature check.}
The complete supplied definition, including source qualifications and the dated literature assessment, was reviewed. The calculation below is an elementary restricted attempt motivated by [S2]; it does not claim novelty or a fresh literature-status determination.

\paragraph{Context and investigated scope.} The complete supplied definition, including source qualifications and the dated literature assessment, was reviewed. The calculation below is an elementary restricted attempt motivated by [S2]; it does not claim novelty or a fresh literature-status determination.

\paragraph{A nonsemisimple sewing-algebra test.} Let $A=\mathbb C[\epsilon]/(\epsilon^2)$ and define $\tau(1)=0$, $\tau(\epsilon)=1$. The bilinear pairing $B(a,b)=\tau(ab)$ has matrix $\left(\begin{smallmatrix}0&1\\1&0\end{smallmatrix}\right)$ in the basis $1,\epsilon$, so is nondegenerate although the algebra is not semisimple. Associativity gives $B(ab,c)=B(a,bc)$; therefore the elementary four-point contractions determined by multiplication and this pairing have the same value $\tau(abcd)$ in the two parenthesizations.

\paragraph{A concrete handle calculation.} The dual basis to $1,\epsilon$ is $\epsilon,1$, so the algebraic handle element is $e=1\epsilon+\epsilon1=2\epsilon$. Hence $e^2=0$, $\tau(e)=2$, and $\tau(e^g)=0$ for $g\geq2$. These are internally consistent finite algebraic contractions. With the involution fixing $\epsilon$, however, the form $\tau(a^*a)$ is not positive: for $a=1-\epsilon$ it equals $-2$. Algebraic gluing identities alone therefore do not produce a positive Hilbert realization.

\paragraph{What this does and does not obstruct.} Logarithmic conformal theories need not be unitary, so lack of this positivity is not a no-go theorem for the source. Conversely, a nondegenerate Frobenius pairing is not a construction of analytic local conformal correlators with prescribed chiral category. One still needs braided/nondegenerate categorical data, bulk objects, all sewing identities and mapping-class-group compatibility, plus analytic field control. The first gap is upgrading such algebraic candidates to the full source realization without silently adding hypotheses to every finite tensor category.

\paragraph{Verification and adversarial audit.}
Checked pairing determinant, associativity contractions, dual basis/handle element and negative norm; did not impose unitarity on logarithmic CFT. Model self-audit only; no independent human proof checking. Prompts A and B were applied to this canonical record and its qualified source scope. Checked pairing determinant, associativity contractions, dual basis/handle element and negative norm; did not impose unitarity on logarithmic CFT. Retained unresolved status and zero original-target progress; No full bulk/correlator construction with prescribed chiral category and all sewing/mapping-class/analytic properties is obtained.

Coordinator checked the dual-number Frobenius pairing, dual basis, handle contractions and negative value of the proposed Hermitian form. This additional check was model-assisted, not independent human verification.

\paragraph{Tools, process and contributions.}
Python 3 was used for file preparation and structural checks; no numerical experiment or theorem prover was used. Contribution prompt A was followed by prompt B. GPT-6 Astra with ultra reasoning drafted the research and performed its stated self-audit; any additional coordinator contribution is identified in the verification and audit. A GPT-6-based Codex coordinator assembled the manuscript and checked repository integration; its exact variant and reasoning configuration were not exposed. No human mathematical contribution, independent human proof check or formal proof certification is supplied. Complete derivations and the consulted source references are included here; private reasoning transcripts are not published.

\paragraph{Final status and first remaining gap.}
\textbf{UNRESOLVED}. No full bulk/correlator construction with prescribed chiral category and all sewing/mapping-class/analytic properties is obtained.

\section{Completion Estimate}
\noindent\textbf{COMPLETION: 0\%}
The calculation is an auxiliary or restricted attempt and establishes no newly verified portion of the full source target. No full bulk/correlator construction with prescribed chiral category and all sewing/mapping-class/analytic properties is obtained.

\subsection{Sources}
\begin{itemize}
\item[\hypertarget{source:30002839:1}{S1}] ULAM.AI and the UnsolvedMath Contributors, UnsolvedMath, record 30002839 (OWR-13673-008). Catalogue curation under CC BY 4.0; definitions and English exposition added locally.
\url{https://huggingface.co/datasets/ulamai/UnsolvedMath}.
\item[\hypertarget{source:30002839:2}{S2}] Subfactors and Conformal Field Theory (2015), Oberwolfach Reports, DOI 10.4171/owr/2015/16, pp. 849--926. Primary report accessed 2026-10-11; relevant PDF pages 32, 33, 34.
\url{https://doi.org/10.4171/owr/2015/16}.
\item[\hypertarget{source:30002839:3}{S3}] Gaberdiel--Kausch, local logarithmic CFT (catalogue-listed context; not independently consulted in this preparation).
\url{https://doi.org/10.1016/S0550-3213(98)00701-9}.
\item[\hypertarget{source:30002839:4}{S4}] Fuchs--Schweigert, non-semisimple correlators (catalogue-listed context; not independently consulted in this preparation).
\url{https://doi.org/10.1016/j.aim.2016.11.020}.
\end{itemize}
\label{end:30002839}
\end{document}
